\documentclass[11pt]{article}
% Shared preamble of the RMC kernel write-ups (% Shared preamble of the RMC kernel write-ups (% Shared preamble of the RMC kernel write-ups (\input{preamble} after \documentclass).
\usepackage[margin=1in]{geometry}
\usepackage{amsmath,amssymb,bm}
\usepackage{graphicx}
\usepackage{booktabs}
\usepackage{hyperref}

\newcommand{\dd}{\,\mathrm{d}}
\newcommand{\Ein}{E}
\newcommand{\Eout}{E'}
\newcommand{\Ecm}{E_{c}}
\newcommand{\mucm}{\mu_{c}}
\newcommand{\mulab}{\mu_0}

% Common closing section: how the verification figures are produced
\newcommand{\verificationmethod}{%
Each figure is produced by \texttt{verification/verify\_laws.py}. For an incident
energy $\Ein$ along $+z$, $10^6$ collisions are sampled with MC/DC's own reaction
sampler (\texttt{sample\_elastic\_scattering}, \texttt{sample\_inelastic\_scattering}
or \texttt{sample\_fission}), and every emitted neutron is histogrammed in
$48 \times 24$ bins of $(\Eout, \mulab)$, with $\mulab$ the lab cosine to $+z$. The
inverted kernel used by RMC is integrated over the same bins (Gauss--Legendre panels;
two-body lines are split where $\mulab^*(\Eout)$ crosses a cosine edge), multiplied by
the number of collisions, and compared with a $\chi^2$ test over bins with at least 5
expected counts (the rest pooled into one bin). The tables list $\chi^2/\mathrm{dof}$,
its $p$-value, the emitted neutrons per collision (sampled vs.\ the integral of the
inverted kernel), and the largest relative change of a bin integral when the
quadrature panels are doubled. Left panels: outgoing-energy marginal; middle:
lab-cosine marginal; right: per-bin $z = (\text{sampled} - \text{inverted}) /
\sqrt{\text{inverted}}$.}
 after \documentclass).
\usepackage[margin=1in]{geometry}
\usepackage{amsmath,amssymb,bm}
\usepackage{graphicx}
\usepackage{booktabs}
\usepackage{hyperref}

\newcommand{\dd}{\,\mathrm{d}}
\newcommand{\Ein}{E}
\newcommand{\Eout}{E'}
\newcommand{\Ecm}{E_{c}}
\newcommand{\mucm}{\mu_{c}}
\newcommand{\mulab}{\mu_0}

% Common closing section: how the verification figures are produced
\newcommand{\verificationmethod}{%
Each figure is produced by \texttt{verification/verify\_laws.py}. For an incident
energy $\Ein$ along $+z$, $10^6$ collisions are sampled with MC/DC's own reaction
sampler (\texttt{sample\_elastic\_scattering}, \texttt{sample\_inelastic\_scattering}
or \texttt{sample\_fission}), and every emitted neutron is histogrammed in
$48 \times 24$ bins of $(\Eout, \mulab)$, with $\mulab$ the lab cosine to $+z$. The
inverted kernel used by RMC is integrated over the same bins (Gauss--Legendre panels;
two-body lines are split where $\mulab^*(\Eout)$ crosses a cosine edge), multiplied by
the number of collisions, and compared with a $\chi^2$ test over bins with at least 5
expected counts (the rest pooled into one bin). The tables list $\chi^2/\mathrm{dof}$,
its $p$-value, the emitted neutrons per collision (sampled vs.\ the integral of the
inverted kernel), and the largest relative change of a bin integral when the
quadrature panels are doubled. Left panels: outgoing-energy marginal; middle:
lab-cosine marginal; right: per-bin $z = (\text{sampled} - \text{inverted}) /
\sqrt{\text{inverted}}$.}
 after \documentclass).
\usepackage[margin=1in]{geometry}
\usepackage{amsmath,amssymb,bm}
\usepackage{graphicx}
\usepackage{booktabs}
\usepackage{hyperref}

\newcommand{\dd}{\,\mathrm{d}}
\newcommand{\Ein}{E}
\newcommand{\Eout}{E'}
\newcommand{\Ecm}{E_{c}}
\newcommand{\mucm}{\mu_{c}}
\newcommand{\mulab}{\mu_0}

% Common closing section: how the verification figures are produced
\newcommand{\verificationmethod}{%
Each figure is produced by \texttt{verification/verify\_laws.py}. For an incident
energy $\Ein$ along $+z$, $10^6$ collisions are sampled with MC/DC's own reaction
sampler (\texttt{sample\_elastic\_scattering}, \texttt{sample\_inelastic\_scattering}
or \texttt{sample\_fission}), and every emitted neutron is histogrammed in
$48 \times 24$ bins of $(\Eout, \mulab)$, with $\mulab$ the lab cosine to $+z$. The
inverted kernel used by RMC is integrated over the same bins (Gauss--Legendre panels;
two-body lines are split where $\mulab^*(\Eout)$ crosses a cosine edge), multiplied by
the number of collisions, and compared with a $\chi^2$ test over bins with at least 5
expected counts (the rest pooled into one bin). The tables list $\chi^2/\mathrm{dof}$,
its $p$-value, the emitted neutrons per collision (sampled vs.\ the integral of the
inverted kernel), and the largest relative change of a bin integral when the
quadrature panels are doubled. Left panels: outgoing-energy marginal; middle:
lab-cosine marginal; right: per-bin $z = (\text{sampled} - \text{inverted}) /
\sqrt{\text{inverted}}$.}

\usepackage{amsthm}

\newtheorem{proposition}{Proposition}
\newcommand{\tpsi}{\tilde\psi}
\newcommand{\teps}{\tilde\epsilon}
\newcommand{\Sbar}{\bar S}
\newcommand{\Sigt}{\Sigma_t}

\title{Linear discontinuous energy trial space for RMC}
\date{}

\begin{document}
\maketitle

\noindent Code: \texttt{energy\_basis="linear"} in \texttt{mcdc.rmc.run};
\texttt{\_add\_outgoing}, \texttt{\_add\_incoming}, \texttt{reaction\_transfer\_row},
\texttt{nuclide\_transfer\_moments} in \texttt{kernel.py};
\texttt{binned\_source}, \texttt{projected\_removal}, \texttt{collision\_residual\_mass},
\texttt{\_abs\_quadratic\_integral}, \texttt{Residual} in \texttt{residual.py};
\texttt{sample\_collision\_edge}, \texttt{in\_scatter\_density},
\texttt{binned\_in\_scatter}, \texttt{correction\_masses},
\texttt{sample\_correction\_integrated}, \texttt{sample\_correction} in
\texttt{source.py}; MC/DC track-length score \texttt{flux-energy-slope}
(\texttt{mcdc/transport/tally/score.py}). Notation as in \texttt{rmc\_math.tex}:
$E'$ incident and $E$ outgoing energy in the transport equations, cells $k$, energy
bins $g$, polar-cosine bins $j$.

\section{Motivation}

RMC stalls when the part of the residual that cannot vanish, $r_\infty = q - L P\psi$,
dominates: at $P\psi$ the exact correction has zero projected mean, but generally
nonzero finite-history variance (\texttt{rmc\_math.tex}, Section~2). That variance
depends on the residual, proposal and transport tally response. With a
piecewise-constant trial space, $\|r_\infty\|$ is set by the in-bin shape of the flux,
$\psi - P\psi = O(\Delta E\,\partial_E\psi)$. A linear trial space reduces this to
$O(\Delta E^2\,\partial_E^2\psi)$ wherever $\psi$ is sufficiently smooth in the bin.
This can reduce the unresolved residual, but does not guarantee a smaller tally
variance or convergence rate. A line generally resolves a narrow resonance dip
poorly; a proposed alternative is the $1/\Sigt(E)$-weighted polynomial basis of
\texttt{self\_shielded\_basis.tex}. An additional $1/E$ factor would be a separate
spectral assumption.

\section{Continuous or discontinuous}

The exponentially convergent Monte Carlo (ECMC) work at Texas A\&M used linear
\emph{discontinuous} (LD) trial spaces in space and angle (Peterson, Morel and Ragusa,
M\&C 2013; Franke, Bruss and Morel, ANS 2013; Bolding and Morel, M\&C 2017), with
h-refinement to maintain convergence after inadequate resolution causes stagnation;
see the published space--direction formulation \cite{peterson}. The appropriate
refinement indicator depends on the algorithm. Vermaak and Morel
\cite{vermaak} showed that \emph{spatial}
discontinuities are a pitfall. The streaming term turns a jump $D$ at face $z_f$ into
the face source $-\mu D\,\delta(z-z_f)$, which, seen from the two adjacent cells, is a
pair of nearly equal sources of opposite sign; their responses cancel and only add
variance. On a 1D slab, piecewise-constant RMC was $6.5\times$ \emph{less} efficient
than SMC, continuous linear $3\times$ more. These are results for their
Sec.~4.1 slab benchmark, not general efficiency factors.

The transport operator differentiates only in $z$: there is no $\partial_E$ (no
continuous slowing down) and, in slab geometry, no $\partial_\mu$. Jumps in $E$ and
$\mu$ therefore create no distributional derivative residual, while the true flux can
have jumps or singularities there
(at a source energy $E_0$, at the elastic edge $\alpha E_0$, at discrete levels, and
at $\mu=0$ at vacuum boundaries). The choice is therefore
\begin{itemize}
  \item \textbf{linear discontinuous in $E$} (this note) and in $\mu$ (planned): local
  projection, block-diagonal mass matrix, no extra residual terms;
  \item \textbf{linear continuous in $z$} (planned): no interior face residual; the
  projection is a tridiagonal solve per $(g,j)$.
\end{itemize}

\section{Basis and projection}

On bin $g=[E_g,E_{g+1}]$ with width $\Delta E_g$, let
\begin{equation}
  x_g(E) = \frac{2E - E_g - E_{g+1}}{\Delta E_g}\in[-1,1],\qquad
  P_0 = 1,\quad P_1(x) = x ,
\end{equation}
and represent
\begin{equation}
  \tpsi(z,E,\mu) = \sum_{a=0}^{p} \tpsi_{kgj,a}\,P_a\big(x_g(E)\big),
  \qquad (z,E,\mu)\in(k,g,j),
\end{equation}
with $p=0$ (constant) or $p=1$ (linear). Since
$\int_g P_aP_b\dd E = \Delta E_g\,\delta_{ab}/(2a+1)$, the $L^2$ projection is local
and diagonal:
\begin{equation}
  (Pf)_{kgj,a} = \frac{2a+1}{h_k\Delta E_g\Delta\mu_j}
  \int_k\!\dd z\int_g\!\dd E\int_j\!\dd\mu\; f(z,E,\mu)\,P_a\big(x_g(E)\big).
  \label{eq:projection}
\end{equation}
$\tpsi_{kgj,0}$ is the bin average and $\tpsi_{kgj,1}$ half the change across the
bin. Errors are reported in the function norm
$\|\tpsi\|^2 = \sum\sum_a \tpsi_{kgj,a}^2/(2a+1)$ (per unit bin volume), which reduces
to the coefficient norm for $p=0$.

\paragraph{Estimating the coefficients.} The energy of a particle is constant along a
track, so a track-length estimator of \eqref{eq:projection} scores
$w\,\ell\,P_a(x_g(E))$. MC/DC's \texttt{flux} score gives $a=0$; the score
\texttt{flux-energy-slope} added for this purpose gives $a=1$ (it requires an
energy filter, whose bin supplies $E_g, E_{g+1}$). Then
\begin{equation}
  \teps_{kgj,a} = (2a+1)\,\frac{\text{tally}_{kgj,a}}{h_k\Delta E_g\Delta\mu_j}.
\end{equation}

\section{Transfer moments}

The in-scatter source of $\tpsi$ projects exactly (up to quadrature) onto the trial
space through the moments
\begin{equation}
  M^r_{a'g'j',\,agj} = \int_{g'}\!\dd E'\,\sigma_r(E')\,P_{a'}\big(x_{g'}(E')\big)
  \int_{j'}\!\dd\mu'\int_g\!\dd E\,P_a\big(x_g(E)\big)\int_j\!\dd\mu\;
  \tau_r(E'\!\to E,\mu'\!\to\mu),
  \label{eq:moments}
\end{equation}
\begin{equation}
  \Sbar_{kgj,a} = \frac{2a+1}{\Delta E_g\Delta\mu_j}
  \sum_{a'g'j'} M_{m_k,\,a'g'j',\,agj}\;\tpsi_{kg'j',a'} = \big(P\mathcal S\tpsi\big)_{kgj,a}.
\end{equation}
Here $M_m=\sum_{r\in m}N_rM^r$, with $N_r$ the atom density of reaction $r$'s
nuclide, as in \texttt{rmc\_math.tex}.
The $a'=a=0$ block is the piecewise-constant moment matrix. Storage grows by
$(p+1)^2$, i.e.\ $4\times$.

\paragraph{Quadrature.} The kernel evaluations are unchanged; only the accumulation
is weighted.
\begin{itemize}
  \item \emph{Outgoing weight.} Every inner quadrature node knows its lab outgoing
  energy: elastic uses the speed--energy map in \texttt{elastic.tex};
  $E = E'\,(A^2+2A\mucm+1)/(A+1)^2$ is its classical limit. Level and continuous
  COM laws use the classical COM-to-lab map; lab and free-gas laws integrate in $E$
  directly. The
  node's contribution $w\,A_{jj'}$ is added to both $a=0$ and, multiplied by
  $x_g(E)$, $a=1$ (\texttt{\_add\_outgoing}). For classical COM kinematics,
  $E(\mucm)$ is affine at fixed $E'$ and $\Ecm$, so the outgoing weight is linear
  in that variable. The full integrand also contains the angular transfer and
  need not be polynomial. MC/DC's elastic energy map is not affine, so neither
  a polynomial-degree argument nor Gaussian exactness applies to that map.
  \item \emph{Incident weight.} The adaptive G7--K15 integral over $E'$ accumulates
  each integrand $f(E')$ into $K_0 \mathrel{+}= w f$ and
  $K_1 \mathrel{+}= w\,x_{g'}(E') f$ (\texttt{\_add\_incoming}); the error estimate and
  acceptance use all blocks together.
\end{itemize}

\paragraph{Checks.} For any reaction that conserves particles inside the grid,
$\sum_{gj}M^r_{1g'j',0gj} = \Delta\mu_{j'}\int_{g'}\sigma_r\,x_{g'}\dd E'$ (the
incident-weighted row sum), and $|M^r_{0g'j',1gj}|\le M^r_{0g'j',0gj}$ (because the
kernel is non-negative and $|x_g|\le 1$). Both are tested on O-16 elastic
(\texttt{test\_elastic\_row\_linear\_basis}).

\paragraph{Cache.} The cache key includes the basis order \texttt{energy\_order};
the moment format version is 3.

\section{Residual}

\subsection{Collision residual}

Inside bin $(k,g,j)$,
\begin{equation}
  r_c(E) = c(E) - \Sigt(E)\,\tpsi(E),\qquad
  c(E) = \sum_a c_{kgj,a}P_a\big(x_g(E)\big),\quad
  c_{kgj,a} = Q_{kgj,a} + \Sbar_{kgj,a} + \bar F_{kgj,a}.
\end{equation}
The external source enters through its own projection $Q_{kgj,a}$; a source uniform
in a bin has $Q_{kgj,1}=0$. For a general source, the full residual also contains
$r_q=q-Pq$. Omitting this remainder solves the projected-source problem.

\paragraph{Exact sampling mass.} $\Sigt$ is linear between the points $E_i$ of the
material's union grid, and $c$ and $\tpsi$ are linear in $E$, so $r_c$ is a
\emph{quadratic} on each interval $[E_i,E_{i+1}]\cap g$. With
$y_0, y_m, y_1$ its values at the ends and the midpoint, and $t\in[0,1]$,
\begin{equation}
  y(t) = y_0 + b\,t + a\,t^2,\qquad a = 2y_0-4y_m+2y_1,\quad b = -3y_0+4y_m-y_1 .
\end{equation}
The roots in $(0,1)$ are found in a cancellation-free form,
$q=-\tfrac12\big(b+s_b\sqrt{b^2-4ay_0}\big)$, $t_1=q/a$,
$t_2=y_0/q$, where $s_b=1$ for $b\ge0$ and $-1$ otherwise. Only real roots in
$(0,1)$ are retained; $q=0$ gives the repeated zero root, and a negligible $a$
uses the linear root $t=-y_0/b$ if $b\ne0$. Between consecutive roots $|y|$
is a quadratic, which Simpson's rule integrates exactly:
\begin{equation}
  \int_{E_i}^{E_{i+1}}|r_c|\dd E = (E_{i+1}-E_i)\sum_{\text{pieces}}
  \frac{t_b-t_a}{6}\Big(|y(t_a)| + 4\,\big|y\big(\tfrac{t_a+t_b}{2}\big)\big| + |y(t_b)|\Big).
\end{equation}
The bin mass is $A_{kgj} = h_k\Delta\mu_j\int_g|r_c|\dd E$
(\texttt{collision\_residual\_mass}).

\paragraph{Sampling.} As before, a bin is drawn with probability $A_{kgj}/A$ and
$(z,E,\mu)$ uniformly inside it, so the weight is
$w = (N/N_c)\,r_c(E)\,h_k\Delta E_g\Delta\mu_j\,A/A_{kgj}$, now with $r_c(E)$
evaluated from both coefficients.

\subsection{Projected removal (collision-only iterations)}

The collision-only phase replaces $\Sigt\tpsi$ by its projection,
\begin{equation}
  \big(P[\Sigt\tpsi]\big)_{kgj,a} = \sum_b R_{m_k,g,ab}\,\tpsi_{kgj,b},\qquad
  R_{g,ab} = \frac{2a+1}{\Delta E_g}\int_g\Sigt(E)\,P_aP_b\dd E ,
\end{equation}
computed exactly by Simpson's rule on the lin-lin grid (the integrand is at most
cubic; \texttt{projected\_removal}). This term is folded into $c$, so the sampled
$r_c = \sum_a\big(c_a - (R\tpsi)_a\big)P_a$ is linear in each bin. For $p=0$,
$R_{g,00}$ is the bin-averaged $\Sigt$ of the piecewise-constant scheme.

\begin{proposition}[Fixed point of the collision-only iteration]
With the projected removal and no scattering correction, the residual lies in the
trial space. In an infinite medium the iteration is then a stochastic fixed-point
iteration for the Galerkin system
\[
  \sum_b R_{g,ab}\,\tpsi_{gj,b} = Q_{gj,a} + \Sbar_{gj,a}[\tpsi]
  \qquad\text{for all } g, j, a ,
\]
which is the discontinuous-Galerkin-in-energy discretization of the slowing-down
equation with exact cross sections and moments (provided $PL^{-1}$ is nonsingular on
the trial space, so that $PL^{-1}r=0$ implies $r=0$). Its residual vanishes at the
fixed point. Mean convergence additionally requires
$\rho(I-PL^{-1}L_h)<1$, and geometric decay of the stochastic error requires
contraction in mean square for the chosen proposal and history count; see
\texttt{rmc\_math.tex}. The zero-residual property alone implies neither condition.
\end{proposition}

For $p=0$ this is the flat-weighted multigroup solution (with the Jensen bias in
resonance bins described in \texttt{rmc\_math.tex}). For $p=1$ the approximation
error can be one order higher in $\Delta E$ for sufficiently smooth fluxes
under a stable Galerkin approximation.

\subsection{Edge residual}

The face jump inherits the energy dependence of $\tpsi$,
$D_{fgj}(E) = D_{fgj,0} + D_{fgj,1}\,x_g(E)$, and $r_e = -\mu D(E)\,\delta(z-z_f)$.
The mass uses the exact mean of $|D|$ over the bin,
\begin{equation}
  \frac1{\Delta E_g}\int_g|D|\dd E =
  \begin{cases} |D_0|, & |D_1|\le|D_0|,\\[2pt]
  \dfrac{D_0^2+D_1^2}{2|D_1|}, & \text{otherwise},\end{cases}
  \qquad B_{fgj} = \Delta E_g\,\overline{|D|}\int_j|\mu|\dd\mu .
\end{equation}
$E$ is sampled uniformly in the bin and $\mu\propto|\mu|$. The density is
$q = (B_{fgj}/B)\,\Delta E_g^{-1}\,|\mu|/\!\int_j|\mu|\dd\mu$, so
\begin{equation}
  w = \frac{N}{N_c}\,\operatorname{sign}(-\mu)\,D(E)\,
  \frac{B\,\Delta E_g\int_j|\mu|\dd\mu}{B_{fgj}} ,
\end{equation}
which reduces to $(N/N_c)\,B\operatorname{sign}(-\mu D)$ for a constant jump.

\subsection{Scattering and fission corrections}

The remainder $r_s = T - \Sbar$ now has a linear $\Sbar(E) = \sum_a\Sbar_aP_a(x_g(E))$,
and the pointwise in-scatter density uses the linear incident flux:
\begin{equation}
  T(E,\mu) = \sum_{g'j'}\int_{g'}\dd E'\;
  \tpsi_{kg'j'}(E')\,U_{j'}(E'\!\to E,\mu),\qquad
  \tpsi_{kg'j'}(E') = \sum_{a'}\tpsi_{kg'j',a'}P_{a'}\big(x_{g'}(E')\big).
\end{equation}
\begin{itemize}
  \item \emph{Integrated sampler} (default): unchanged apart from these two
  evaluations; the pilot masses use $|T-\Sbar(E)|$ at $2\times2$ Gauss points plus the
  floor $0.1\,h\Delta E\Delta\mu\sum_a|\Sbar_a|$. For exact moments,
  $\int_g\dd E\int_j\dd\mu\,P_a(x_g(E))r_s=0$ for $a=0,1$, so $r_s$ carries
  the shape beyond linear. Signed cancellation can make both coefficients and
  all pilot samples zero while $r_s\ne0$ elsewhere: the floor is not a support
  guarantee. An independent proposal positive on the possible support is needed
  if that support is not otherwise established. Finite quadrature qualifies
  both the moment identities and unbiasedness.
  \item \emph{Pointwise sampler}: the subtracted density per unit $E'$ is
  \begin{equation}
    \tilde s(E',E) = \frac{1}{\Delta E_g\Delta\mu_j\Delta E_{g'}}
    \sum_{a'j'a}(2a+1)\,P_a\big(x_g(E)\big)\,M_{a'g'j',agj}\,\tpsi_{kg'j',a'} ,
  \end{equation}
  whose $E'$ integral over $g'$, summed over $g'$, is $\Sbar(E)$. The defensive proposal
  must stay positive wherever the integrand can be nonzero: $q_S$ selects bins
  $\propto h\sum_a(2a+1)\big|\sum_{a'j'}M\tpsi\big|$ (an upper bound of $|\tilde s|$
  times the bin volume), and $q_T$ selects $j'$ $\propto\sum_a|\tpsi_{kg'j',a}|\Delta\mu_{j'}$,
  which is positive wherever $\tpsi_{kg'j'}(E')\ne 0$.
\end{itemize}

\section{Verification}

The tests and numerical results below describe existing verification artifacts;
they have not been rerun for this mathematical revision.

\paragraph{Unit tests.}
\begin{itemize}
  \item \texttt{test\_abs\_quadratic\_integral}, \texttt{test\_collision\_residual\_mass}:
  the exact $|r_c|$ integrals against fine trapezoid rules (constant and linear
  coefficients, rel.\ $10^{-8}$).
  \item \texttt{test\_projected\_removal}, \texttt{test\_binned\_source}: $R$ and $\Sbar_a$
  against direct sums and integrals.
  \item \texttt{test\_collision\_edge\_unbiased[P=2]}: the mean weight of the collision
  and edge particles per bin equals the exact bin integral of $r_c$ and $r_e$ (5$\sigma$).
  \item \texttt{test\_elastic\_row\_linear\_basis}: the moment checks above.
  \item \texttt{test\_analytic\_slowing\_down\_A1\_linear}: end to end, below.
\end{itemize}

\paragraph{Analytic slowing down.} $A=1$, $\sigma_s=\sigma_c$ constant, isotropic
scattering, unit source at $E_0=101$\,eV smeared over $[E_0-0.1, E_0]$; below the
source, the classical line-source limit is
$\phi(E) = \tfrac{c}{\Sigt E_0}(E_0/E)^c$ with $c=1/2$.
For the stated uniform unit source on $[E_a,E_b]=[100.9,101]$\,eV, its exact
classical expression below $E_a$ is
$\phi(E)=(E_b^c-E_a^c)/[\Sigt(E_b-E_a)E^c]$.
The scaling and sampler-kinematics qualifications in \texttt{rmc\_math.tex} apply.
There are 20 log-spaced
bins over $[1, 100.9]$\,eV, 2 polar bins, and 500 histories per bin per iteration.

\begin{center}
\begin{tabular}{lcc}
\toprule
 & constant & linear \\
\midrule
Collision-only fixed point, max.\ rel.\ error of bin averages & $3\times10^{-3}$ & $2.5\times10^{-4}$ \\
\quad max.\ rel.\ error of $P_1$ coefficients & -- & $9\times10^{-4}$ \\
Full-residual iterations, $\|\teps\|$ before the stall
  & $3.5,\ 1.3\!\times\!10^{-2},\ 7.8\!\times\!10^{-4}$
  & $3.5,\ 0.19,\ 1.2\!\times\!10^{-2},\ 5.3\!\times\!10^{-4},\ 8.5\!\times\!10^{-5}$ \\
\quad stall level of $\|\teps\|$ & $\sim10^{-3}$ & $\sim7\times10^{-5}$ \\
\quad rms rel.\ error of bin averages, iterations 0--3
  & $6.3,\ 1.0,\ 1.1,\ 0.84\ (\times10^{-2})$
  & $6.3,\ 1.4,\ 0.15,\ 0.035\ (\times10^{-2})$ \\
\quad rms rel.\ error at the stall & $0.7$--$1\times10^{-2}$ & $\approx5\times10^{-4}$ \\
\bottomrule
\end{tabular}
\end{center}

\noindent The full-residual iteration (collision, edge and integrated scattering
correction in every iteration) converges for four iterations with the linear basis
against two with the constant basis, and stalls about $15\times$ lower in
$\|\teps\|$ and $15$--$20\times$ lower in the true error of the bin averages, consistent
with $\|r_\infty\|$ dropping from $O(\Delta E)$ to $O(\Delta E^2)$. The collision-only
iteration reduces $\|\teps\|$ by about $10\times$ per iteration (the constant basis by
about $20\times$).

\noindent These are previously recorded results, not new verification of the
corrected mathematical statements. In particular, a source projection and the
classical reference model must be distinguished from the exact physical residual
and the implemented elastic sampler.

\begin{thebibliography}{9}
\bibitem{peterson} J.~R.~Peterson, J.~E.~Morel and J.~C.~Ragusa,
  \emph{Residual Monte Carlo for the One-Dimensional Particle Transport Equation},
  SIAM J.\ Sci.\ Comput.\ \textbf{38}(6), B941--B961 (2016),
  \href{https://doi.org/10.1137/16M1057619}{doi:10.1137/16M1057619}.
\bibitem{vermaak} J.~I.~C.~Vermaak and J.~E.~Morel,
  \emph{Transport Error Estimation Using Residual Monte Carlo},
  J.\ Comput.\ Phys.\ (2022),
  \href{https://doi.org/10.1016/j.jcp.2022.111306}{doi:10.1016/j.jcp.2022.111306},
  Secs.~3.4 and 4.1.
\end{thebibliography}

\end{document}
